\documentclass[]{../../../../tex/classes/styledArticle}
\usepackage{../../../../tex/packages/hebrewSupport}

\author{שחר פרץ}
\title{\textit{היסטוריה 42}}
\date{12 לנובמבר 2025}
\begin{document}
	\maketitle
	\section*{מלחמת העצמאות}
	זו מלחמת קיום, לכל אורכה. נפריד בין שני דברים: 
	\begin{itemize}
		\item הפלסטינים שלא מכירים בהחלטה, ואומרים שהם יעשו הכל כדי לעצור אותה. ואכן הם יצאו למלחמה יום אחרי ההחלטה. 
		\item מדינות ערב שגם הן לא הכירו בהחלטה, אך אמרו שרק אם תוקם מדינה יפלשו. ואכן, המלחמה בצבאות ערב הסדירים החלה אחרי ההכרזה עם קום המדינה. 
	\end{itemize}
	השלב הראשון במלחמה, אינו ממש מלחמה. ''יותר כמו אינטיפאדה``. צלפים בערים מעורבות, אוטובוסים מתפוצצים וכו'. אין אסטרטגיה, אלא רק 4 חודשים לתוך המלחמה. החלק הראשון \textbf{הוא מלחמת אזרחים}. בתוך שני החלקים האלו, לפני ואחרי הכרזת המדינה, אנחנו נחלק את המלחמה לשני תתי חלקים, בהתאם למצב של היישוב במלחמה. במקפיים בעמוד הראשון של הדף ששרית חלקה יש את האירועים שבאמת רלוונטיים לנו למבחן. 
	העובדה שהמלחמה מלחמת קיום (או שאנחנו פה, או שאנחנו לא. בלי פה. אין לאן ללכת) השפיעה על הלחימה. לאורך כל המלחמה, אין הרדה בין חזית לעורף. אין אוכלוסיה אזרחית שלא במידה כזו או אחרת נמצאה בתוך המלחמה, והלחימה התרחשה גם בתוך היישובים. 
	
	ביישוב היהודי יש כ־600 אלף אנשים, לעומת הערבי עם 1.2 מיליון. \textit{זה לא מייצג את יחסי הכוחות}. מבחינה דמוגרפית, האוכלוסיה הערבית בנויה אחרת – יש בה הרבה יותר ילדים, הרבה יותר מבוגרים, וחצי ממנה הוא נשים (כצפוי). האוכלוסיה אצלנו עם ילודה יותר נמוכה (היא מתפתחת), זקנים לא בהקפים שיש באוכלוסיה הערבית, ואצלנו נשים נלחמו. כלומר, מבחינת הפוטנציאל בכוח לוחם, המצב היה לטובתנו: היו יותר אנשים בטווח גילאים רלוונטי + נשים נלחמו. סה''כ המספרים לא באמת היו ''הבדלים בשמיים``, והכוח הלוחם מהצד היהודי היה הרבה יותר מיומן, בדגש על הפלמ''ח – הוא מאורגן בגדודים, מתאמנים, יש פיקוד, התארגנות, וכו'. תוך כדי הגיעו גם אנשים מאירופה שנלחמו במלחמת העולם השנייה, כגון פרטיזנים ואנשים שנלחמו בצבאות בעלות הברית. מאוחר יותר (בחלק השני, לא במלחמת האזרחים) יש מח''ל – מתנדבי חוץ־לארץ. בד''כ קצינים שהגיעו מחו''ל כדי לסייע בלחימה. אפשר להוסיף לזה את העלייה מקפריסין (אומנם כשהבריטים היו כאן לא היה אפשר לצייד אותם בנשק, אבל הם כן עברו אימונים ע''י אנשי ההגנה). 
	
	
	\subsection*{מלחמת האזרחים}
	את הלחימה במלחמת האזרחים אפיינה לחימה שנקראת ''פזעה`` – קריאה. האוכלוסיה הייתה גדולה, אך לא היה צבא סדיר. ההמון הוא זה שנלחם. כשהמון מתנפל על סיירה של רכבים, אנשים, או יישוב, גם אם ההמון לא מיומן או אפילו בלי נשק חם – יחסי הכוחות לטובתו. 
	
	\subsubsection*{השלב הראשון}
	השלב הראשון מתרחש כשאין כאן מדינה. השלטון הבריטי (שהקים את היציאה שלו, היה אמור לצאת באוגוסט ולא במאי). האינטרס של הבריטים היה לצאת בלי לפגוע בעצמם, ובהתאם לזה ניהלו את המדיניות שלהם. הבריטים הם בעד הבריטים, ומכורח היותם עדיין פה, היו עדיין מגבלות על נשק והעפלה, שהעיקו על היישוב היהודי לנהל מלחמה (עקרונית עדיין היה אסור לקנות ולהכניס נשק, כי אנחנו לא מדינה ריבונית. זה כמו למכור נשק למיליציות. בפועל נקנה נשק). 
	
	נתבונן במפת החלוקה. אומנם, השטח שלנו יותר גדול, אבל אין רצף טריטוריאלי. המעברים בין החלקים המקוטעים האלו, במקרים רבים, היו בלתי אפשריים. גם בתוך השטח עצמו היו יישובים מבודדים, בגלל שהערבים שלטו על הדרכים. הטופוגרפיה של המדינה הפלסטינית יותר טובה, עיו''ש יותר גבוהים משאר השטח. איזור עכו וצפונה, היו הגליל המערבי (גם גבוהה) – לא בשטחנו. יישובים יהודים כמו ניצנים נמצאו בשטח המדינה הפלסטינית (רובם היו בשטח המדינה היהודית), רובם בנגב ובגליל המערבי. 
	
	ישנן ערים מעורבות – ירושלים, יפו, חיפה, צרפת, עכו, טבריה, וכו'. בערים הללו – מלחמת אזרחים. היו צלפים. יש סיפורים לא מעטים על אנשים שישבו בבית, וכדור טועה נכנס והרג אותם. 
	
	בחודש מרץ, יש כמה שיירות גדולות שנתקלות, ושם היו נפגעים רבים. עד כאן זהו השלב הראשון במלחמת האזרחים. כאן הייתה נקודת מפנה. 
	
	נתבונן באירוע לדוגמה – שיירת הל''ה. זוהי שיירה רגלית (כל השאר, היו שיירות של רכבים, בד''כ 7-10 רכבים ולעיתים אף עשרות). הם יצאו מהרטוב (ליד בית־השמש) לכיוון גוש עציון, ישובים מבודדים שהיו צריכים אוכל וציוד. הם יצאו מאוחר יחסית (וזו הטענה המרכזית כנגד יציאת השיירה – לא היה שום סיכוי בעולם שהם יגיעו לפני הזריחה) והיא מתגלה בדרך (יש וויכוח ע''י מי). הם קיבלו החלטה לא לפגוע במי שגילה אותם, ואלו חזרו לכפר לקרוא למקומיים (חלקם עם נשק, חלקם סתם כפריים לא לוחמים). מתנהל שם קרב מהבוקר עד הלילה, שבסופו כל 35 אנשי השיירה נהרגו. האותה התקופה הערבים התעללו בגופות בצורה זוועתית. בעזרת הבריטים קברו את הגופות מאוחר יותר. 
	
	זו דוגמה: יישובים מבודדים, שליטה של הערבים על הדרכים, ושיירות שמנסות להגיע. באופן דומה בחודש מרץ 1947 נהרגו 80 הרוגים – שיירת חולדה, שיירת יחיעם ושיירת דניאל כולן נתקלו בדרך לירושלים. 
	
	לפני קום המדינה, הייתה תפיסה שלא נוטשים יישובים. הסיבה – כדי לשמור על הגבול. אותה התפיסה של ההתיישבות שנגררה קדימה, איפה שיהיה יישוב שם יעבור הגבול. 
	
	בשלב מסוים האמריקאים נסוגים מהתמיכה בהצעת החלוקה, בגלל כל הבלגן שקרה כאן. הם העלו הצעה שיהיה פה משטר נאמנות של האו''ם ''עד שאפשר שיהיה משהו``. הנהגת היישוב מבינה, שאם לא ישתנה פה משהו, ארצות־הברית גם תעביר את ההחלטה הזו באו''ם, ועוד פעם יהיה כאן גוף לכאורה בינ''ל שישלוט כאן. 
	
	\subsubsection*{השלב השני}
	ב־1 באפריל 1948 התחילה התקופה השנייה. בתקופה הזו מתחילה תוכנית ד' (אחרי תוכניות א', ב' וג'). הרעיון: ליזום פעולות, \textbf{מתוכננות}, עם ממש אסטרטגיה ויעדים ברורים מה אנחנו רוצים להשיג, עד הכרזת המדינה. זה ברור שצבאות ערב יפלשו, ולכל הפחות בהתחלה נהיה בנחיתות (אם לא צבאית, מבחינת נשק) ואנחנו צריכים להגיע לשם עם כמה שיותר הישגים. ההבנות: 
	\begin{itemize}
		\item בראש ובראשונה, אנחנו צריכים רצף טריטוריאלי. בשלב הזה היו עוד יישובים שהתנתקו לגמרי. 
		\item צריך לדאוג שהדרכים לא יהיו בשליטת הערבים. 
		\item צריך לדאוג לערים המעורבות. 
	\end{itemize}
	המבצע הראשון שיצא לדרך בתפיסה הזו, נקרא ''מבצע נחשון`` (הוא חלק מתוכנית ד' אבל לא חלק מתוכנית הלימודים). בניגוד להמלצה של מפקדי הצבא, בן־גוריון קיבל את ההחלטה להעביר 1500 חיילים לטובת פתיחת הדרך לירושלים. מבחינתם – מה קורה עם דגניה וניצנים? בן־גוריון מתעקש, כי הוא רואה בירושלים חשיבות אסטרטגית מבחינה לאומית. ירושלים באותה התקופה במצור, ברמה שיש שם רעב. אין אוכל בעיר, שיירות לא מצליחות להגיע. לא נכנס כרגע למבצע נחשון. במבצע נחשון נפתחה הדרך לירושלים באופן זמני בלבד. הצליחו להעלות לשם שיירות של מאות משאיות, עם ציוד, תחמושת ולוחמים. מאוחר יותר קרב נוסף בו הפסדנו, והדרך לירושלים נסגרה. קצת יותר מאוחר נפתחה דרך בורמה (זמן קצר אחרי הכרזת המדינה) ואז העבירו דרך ואנשים דרך דרך בורמה. 
	
	קרב רמת יוחנן, שעליו נצטרך לקרוא לבד, הוא הקרב הכי חשוב ביצירת מערכת היחסים ביננו לבין הדרוזים. עד הקרב הזה, הדרוזים לא היו בצד שלנו. זה קרב שהתנהל מול הדרוזים מסוריה בעיקר, וגם דרוזים שחיים בארץ, ובסופו נוצרה הברית בין הדרוזים ליישוב היהודי. החל מנקודת המפנה הזו, הדרוזים הם חלק מהכוח הלוחם, ולאחר מכן חלק מצה''ל, והם ממש השתתפו בקרבות במהלך המלחמה. 
	
	בסיום תוכנית ד' לא הושגו כל ההישגים, אבל נכבשו הערים המעורבות (טבריה ראשונה, ואז חיפה ויפו), חלק מהדרכים נפתחו (ירושלים כאמור, נפתח ונסגר), והצליחו להגיע ולעבות חלק מהיישובים המבודדים. בתקופה של תוכנית ד' מתחילה בעיית הפליטים הפלסטינים. פרשת דרייך־סין (או משהו כזה) גרמה לפחד גדול באוכלוסיה הפלסטינית ולפחד גדול שהאירוע יחזור על עצמו. בתקופות יותר מאוחרות, היישוב היהודי גם גירש פלסטינים, אך רובה הייתה בריחה מסיבות שונות (פחד, בשלב יותר מאוחר עידוד של המנהיגים שלהם שיצאו להלחם ואז ינצחו ויחזרו, וכו'). 
	
	\subsection*{לחימה בצבאות סדירים}
	\subsubsection*{הכרזת המדינה}
	באופן יחסי, היישוב היה במצב טוב. היו הישגים, לא ב־100\%, אבל היו הישגים. עדיין, יש התלבטות מאוד גדולה האם להכריז על המדינה בנקודות הזמן הזו. מתנהל שם דיון ''שאני הייתי מאוד רוצה להראות לכם משהו אבל אני לא יכולה [אין מקרן]!``. לא ברור שנערכה הצבעה, אבל דיונים ארוכים בהחלט היו. 
	\begin{itemize}
		\item מצד אחד, היו כאלו שטענו שאי אפשר להכריז עכשיו – זה סיכון גדול מדי. 
		\begin{itemize}
			\item היה ידוע וברור שברגע שנכריז על מדינה, 5 צבאות פולשים לארץ. 
			\item ארה''ב שאיימה להקפיא את הסיוג של יהודי ארה''ב ליישוב אם תתקיים הכרזה על המדינה. היא 
			\item חשש מאברגו – האו''ם ביקש שביתת נשק לשלושה חודשים, ולמדינות ערב היה יתרון צבאי מבחינת נשק. 
			\item גוש עציון נכבש יום לפני הההכרזה ע''י הליגיון הירדני, והוא הצבא החזק ביותר שנלחם נגד מדינת ישראל. ירדן – זה צבא בריטי לכל דבר. מאומן ומחומש ע''י הבריטים. 
			\item ההערכה הייתה שיש לנו סיכוי של 50\%-50\% לשרוד. 
		\end{itemize}
		\item מצד שני, היו כאלו שטענו שחייבים להכריז עכשיו. 
		\begin{itemize}
			\item יש הבנה שאין וואקום, אם הבריטים יצאו מפה ולא נכריז על המדינה, משהו יכנס לפה. אולי שלטון בין־לאומי, אולי מדינות ערב. 
			\item היה חשש שתהיה כאן החמצה לדורות, זו סיטואציה שלא תחזור על עצמה, והיא מוגבלת בזמן. ארצות הברית ''כבר התחילה לעשות קולות של חריקות``. 
			\item מדינה ריבונית יכולה לגייס גיוס חובה. היו משתמטים גם ב־48. לא כולם רצו לקרב ואמרו בואו נקריב את עצמנו. היו כאלו שפרנסה לא אפשרה להם להשתתף. כאשר מכריזים על מדינה, יש אפשרות להכריח אנשים להתגייס. ואכן, פחות משבועיים אחרי הקמת המדינה, הוקם צה''ל עם חוק הגיוס. 
			\item בחוק הבין־לאומי אנחנו לא יכולים לרכוש נשק אלא אם אנחנו מדינה (למרות שכבר בתוכנית ד' היו מבצעים שהביאו לפה נשק, אחד מהם נקרא מבצע חסידה. הנשקים ששמשו במבצע נחשון הוא נשק שהגיע מצ'כוסלובקיה)
			\item עלייה המונית שתהיה כוח לוחם ובונה. 
		\end{itemize}
	\end{itemize}
	בן גוריון העריך, ובצדק, שלעובדה שאנו במלחמת קיום יש משמעות בלחימה ובמורל. יש להם מטרה בזמן הלחימה, זאת בניגוד לצבאות ערב, שם הם נלחמים על דבר שהם אפילו לא יודעים מה הוא. 
	
	בן־גוריון החליט להקים את צה''ל ולפרק את המחתרות. פרשת אלטלנה שאתם כנראה מכירים היא חלק מהנושא הזה. 
	
	\subsubsection*{יחסי הכוחות}
	נתבונן בטבלה בעמוד 97. נבחין שהכוח הצבאי היה בעיקרו לטובת צבאות ערב, בעיקר בנשק כבד – ספינות, שיא, מטוסים, טנקים (לנו היה טנק אחד בלי תותח, להם היו 40). עם זאת, מבחינת כוח אדם לא היה היה חסרון משמעותי עבורנו. 
	
	הכוח שלנו כבר חצי שנה במלחמה, זאת בניגוד לכוחות הרעננים של צבאות ערב. הנשק שלנו היה גם מיושן יותר. אבל!
	\begin{itemize}
		\item צבאות מדינות ערב, לא ממש היה ברור להם על מה הם נלחמים. הבחור האקראי מהנילוס שהציעו לו כסף כדי להלחם בצבא המצרי, בחלקם הגדול לא היה מאומן, או לא ידע למה הוא נלחם. אנשים בצד שלנו היו מוכנים למות במלחמה, החייל המצרי האקראי לא. 
		\item ההנהגה של מדינות ערב הייתה מפוצלת. כל אחת הייתה בעלת האינטרסים שלה. הם נכנסו בהצהרה שהן באות לעזור לפלסטינים, זה ''שטויות במיץ עגבניות`` ''נא לא להקליד את זה עכשיו, שאמרתי שטויות במיץ עגבניות``. היו להם אינטרסים מדיניים משלהם (פרטים תקראו בבית). כל אחת רצתה לכבוש חלק אחר מישראל. 
		\item קווי אספקה ארוכים. ראינו מה זה עשה בסטלינגרד (וולגרד היום). 
		\item חוסר בנסיון קרבי. החיילים ברוב צבאות מדינות ערב, לא היו חיילים שנלחמו, וחלק גדול מהם לא היה גם מאומן יותר מדי. 
	\end{itemize}
	לטובתנו: 
	\begin{itemize}
		\item מוטיבציה גבוהה. 
		\item פיקוד מרכזי אחיד על רוב הכוחות. יש כאן מדינה שקובעת מדיניות, שהצבא מבצע. אומנם, הצבא הלא גדול שלנו פוצל בין כמה חזיתות, אבל הייתה כאן חשיבה אסטרטגית ותיעדוף. 
		\item המון נסיון צבאי. 
		\item חוק גיוס חובה. 
		\item העולים החדשים, שנקראו גח''ל (גיוס חוץ לארץ)
		\item משלוחי נשק גדולים, מצ'כוסלובקיה בעיקר. 
	\end{itemize}
	לרעתנו: 
	\begin{itemize}
		\item הצבא שלנו עדיין בשלבי הקמה. הפקודה יצאה ב־1 ביוני, כשבועיים אחרי הקמת המדינה. מצד אחד קיים כאן צבא (ההגנה והפלמ''ח). מצד שני, זה תהליך שלא הולך חלק. זה לא שהאצ''ל והלח''י ששו ושמחו להתפרק, וגם לא הפלמ''ח. 
		\item מחסור חמור בנשק כבד ובמשאבים כלכליים. 
		\item אנחנו נלחמים בעורף. 
	\end{itemize}
	בתקופה הזאת של המלחמה, ממאי 48 עד יולי 49 (שנה וחודשיים), מגיעה לכאן עלייה המונית. המדינה לא רק במלחמה, היא גם קולטת עלייה. מבחינה כלכלית, זה אתגר מטורף, שאין אף מדינה בעולם שהייתה צריכה להתמודד איתו. כמות העולים יותר גבוהה (משלב מסויים) מכמות הקולטים; הגיעו מיליון עולים, לעומת 650K שהיו כאן קודם לכן. חלקם אולי התחיילו ונכנסו לצה''ל, אבל יש גם ילדים וכו' שלא. 
	
	\subsubsection*{החלק הראשון של הלחימה}
	החודש הראשון הוא חודש של בלימה. מבחינת צבאות ערב, הם התכוננו למלחמת בזק של שבוע, שבועיים או חודש, ולסיים איתנו. זה לא קרה, והם לא היו מוכנים לזה. בשבוע־שבועיים הראשונים היו הצלחות רבות לצבאות ערב – למעלה מאלף הרוגים לכוחתינו. בשלב הזה עוד היו הרבה יישובים מבודדים. 
	
	אירוע לדוגמה: הקרב על ניצנים. ניצנים נמצא ליד ניצנה, הוקם ב־1943. צה''ל החליט כמה ימים לאחר הכרזת המדינה על ''מבצע תינוק`` – פינוי מיישובים מבודדים, בכללם ניצנים, אוכלוסיה בלתי לוחמת וילדים. היו אמהות שהחליטו להישאר. בתחילת יוני 1948 נפלה ניצנים, והלחומים שהיו ביישוב החליטו, בניגוד לרוח שהייתה אז בארץ, לנטוש את היישוב (היה שם סיפור אם כן או לא שלחו כוחות גבעתי). המצרים ניסו כמה פעמים לכבוש את היישוב לפני ההצלחה. הלוחמים, שנכנעו לצבא המצרי, גרר סיפור שלם לאורך הרבה שנים. רוח התקופה הייתה שלא נוטשים יישובים. היה שם סיפור של מספר רב של שנים, שכלל עלבון. אבא קובנר מדבר על ניצנים כבגידה (כמורל לחיילי גבעתי). לאורך הרבה שנים ''לא בדיוק התנהגו אליהם יפה``. כמשתמשים בזה באירוע קרב לדוגמה, צריך לדבר על יישובים מבודדים, יחסי כוחות לטובת צבאות ערב וכוחות רעננים של צבאות ערב לעומת הכוח הישראלי. 
	
	\subsubsection*{החלק השני של הלחימה}
	זה לא היה דבר חד ומהיר, כמו בתקופה השנייה של מלחמת האזרחים. ''אתה פספסת את כל מלחמת העצמאות. אני סיימתי אותה כבר``. ''אוי``. 
	
	אחרי חודש של לחימה נגד צבאות ערב (אנחנו מדברים בשלב הזה על צבאות, ועל האו''ם שמנסה לגשר בין \textit{מדינות} – אלו לא כוחות בלתי סדירים), הצליחו להגיע להפסקת אש של חודש. בחודש הזה, צה''ל מתארגן. מתכננן תוכניות, מביא ספינות נשק ומפיל את חלקן, ויש כאן תכנון ארוך לגבי איפה לכבוש, איפה לרכז כוחות, באיזה כיוון, יעדים ספציפיים, וכו'. יש כאן יוזמה התקפית של צבא מאורגן, לא רק ''לזרום עם המלחמה``. צבאות ערב פחות נצלו את ההפוגה. 
	
	בסיום במלחמה יש תקופה שנקראת ''קרבות עשרת הימים`` (שבמפתיע ארכו עשרה ימים) והצליחו להגיע בהם להישגים לא מבוטלים. כבשו את רמלה ולוד, שם היה גירוש במוצהר של הפלסטינים. במבצע דני (שבהתחלה קראו לא מבצע לרלר), ע''ש דני מס מהל''ה. הכיבוש של רמלה ולוד היה פתיחת דרך לקראת כיבוש הדרום (לא סתם, יש כאן תכנון). לאחר סיום אותם הקרבות הוכרזה הפסקה נוספת, שלא מוגבלת בזמן. היו שיחות להסכמי שביתת נשק. בתוך ההפוגה, מדינת ישראל יוזמת מבצעים גדולים: יואב, חורב, עובדה (כיבוש אילת), וחירם (בגליל העליון). המבצעים גם נועדו לכבוש שטח, וגם להשיג הישגים לקראת שיחות הסכמי שביתת נשק. לדוגמה, במבצע חירם, נכנסו לתוך לבנון וכבשו 14 כפרים – ברור שאף אחד לא רצה לכבוש את לבנון, אבל זה נתן אפשרות למשא ומתן. הדוגמה היא מבצע עובדה, בו לא נורתה ירייה אחת – פשוט הגענו לאילת. 
	
	בעקרון סיימנו את המלחמה. או לכל הפחות את הסקירה של המלחמה. לא דיברנו על בעיית הפליטים ועל ההסכמי שביתת הנשק. יכול להיות שתוקדש בהמשך חצי שעה כדי לדבר על דברים. מומלץ בחום לראות את הדברים האקראיים ששרית תשלח לראות. השיעור הזה הוא רק מבט על, ומכאן ואילך נצטרך להשתמש בספר ובסיכומים. אם אתם רוצים לעשות לעצמכם הנחה, אתם יכולים ללמוד עד ההכרזה כולל (ואז לא לבחור את השאלות של החלק השני. יכול להיות שזה יגרור אותנו להיות מוכרחים לבחור שאלה אחת בלבד). 
	
	
	
	
	
	\ndoc
\end{document}